\setlength{\parindent}{0pt}
\normalsize

Robotic insertion is a common but challenging assembly task because small deviations in position, orientation, tolerance, contact geometry, or force settings can determine whether an insertion succeeds or fails. While machine-learning models can classify robotic time-series data, global classification metrics alone provide limited insight into which physical conditions (experimental and deviational conditions) impact physical insertion success and which conditions are difficult for a model to classify. This thesis therefore constructs and analyzes a real-world robotic insertion dataset that connects multivariate robot time-series data with detailed metadata.

The dataset consists of 8100 peg-in-hole insertion attempts recorded with a UR5e robot under varying rod types, tolerance levels, gripper configurations, recording batches, position offsets, tilt levels, and force settings. Each sample contains 48 synchronized raw robot signals directly from its internal sensors, a binary success label, a failure reason, and metadata describing the corresponding insertion condition. An adapted PatchTST (Patch Time Series Transformer) model is used for insertion-success classification, and the dataset is analyzed through a best channel study, benchmark experiments, and two deviation analyses.

The results show that insertion success is best classified using diverse signal groups, especially force/torque, TCP pose, joint velocity, and current consumption. Classification performance is strong under in-distribution conditions, but decreases for rectangular contact geometries and recording-batch-level shifts. The metadata analysis shows that tilt and force interaction are the strongest factors affecting physical success. It further shows that physical insertion success and classification performance are related but not identical. The thesis demonstrates that combining time-series classification with metadata-based analysis can identify which insertion conditions are impactful to physical success, difficult to classify, and where these two perspectives overlap and differ.
